\documentclass[11pt]{article}
\usepackage{mathblatt}


\blattfuss{Lineare}{Lösungen}
\begin{document}
\noindent{\large\bfseries Lineare \textperiodcentered{} Lösungen}\par\medskip
\begin{pfloesung}{}
\lz{1.}{$A(3|2)$}{}{}
\end{pfloesung}
\begin{pfloesung}{}
\lz{2.}{$9$}{$g(4) = 8 + 1$}{}
\end{pfloesung}
\begin{pfloesung}{}
\lz{3.}{$x^2 + 8x + 16$}{}{}
\end{pfloesung}
\begin{pfloesung}{}
\lz{4.}{$4$}{$x$}{}
\end{pfloesung}
\begin{pfloesung}{}
\lz{5.}{$S(2|3)$}{$x + 1 = -x + 5$, $x = 2$}{}
\end{pfloesung}
\begin{pfloesung}{}
\lz{6.}{$8$}{}{}
\end{pfloesung}
\begin{pfloesung}{}
\lz{7.}{$x_1 = 6$, $x_2 = -6$}{}{}
\end{pfloesung}
\begin{pfloesung}{}
\lz{8.}{$-5$}{$m = 3$, $n$}{}
\end{pfloesung}
\begin{pfloesung}{}
\lz{9.}{steigt}{m ist positiv}{}
\end{pfloesung}
\begin{pfloesung}{{\small\color{mbgrau}P10 2019 · 2b}}
\lz{10.}{Anstieg $-$2: f; parallel zur x-Achse: h}{Steigungsdreieck an f: 1 nach rechts, 2 nach unten}{}
\end{pfloesung}
\begin{pfloesung}{{\small\color{mbgrau}P10 2019 · 2c}}
\lz{11.}{f: y = $-$2x + 2; g: y = x + 2; h: y = $-$2}{y-Achsenabschnitt 2 bei f und g}{}
\end{pfloesung}
\begin{pfloesung}{}
\lz{12.}{$\tfrac{1}{2}$}{$m$; 2 nach rechts, 1 nach oben}{}
\end{pfloesung}
\begin{pfloesung}{{\small\color{mbgrau}P10 2019 · 2a}}
\lz{13.}{x = $-$2}{}{}
\end{pfloesung}
\begin{pfloesung}{}
\lz{14.}{$3$}{$n$}{}
\end{pfloesung}
\begin{pfloesung}{}
\lz{15.}{$2$}{$m$; 1 nach rechts, 2 nach oben}{}
\end{pfloesung}
\begin{pfloesung}{}
\lz{16.}{$-2$}{$m$; 1 nach rechts, 2 nach unten}{}
\end{pfloesung}
\begin{pfloesung}{{\small\color{mbgrau}P10 2024 · 1i}}
\lz{17.}{Punkt (0|$-$1) markiert}{n = $-$1}{}
\end{pfloesung}
\begin{pfloesung}{}
\lz{18.}{$3x - 4$}{$f(x)$}{}
\end{pfloesung}
\begin{pfloesung}{{\small\color{mbgrau}P10 2016 · 1c}}
\lz{19.}{f}{}{}
\end{pfloesung}
\begin{pfloesung}{{\small\color{mbgrau}P10 2025 · 1f}}
\lz{20.}{erste Abbildung (fallende Gerade)}{Gerade ohne Knick}{}
\end{pfloesung}
\begin{pfloesung}{}
\lz{21.}{$3$}{3 nach oben, also $m$}{}
\end{pfloesung}
\begin{pfloesung}{}
\lz{22.}{$-2$}{2 nach unten, also $m$}{}
\end{pfloesung}
\begin{pfloesung}{}
\lz{23.}{$-5$, $-2$, $1$, $4$}{}{}
\end{pfloesung}
\begin{pfloesung}{}
\lz{24.}{Start $(0|4)$, 1 nach rechts, 3 nach unten: Gerade durch $(0|4)$ und $(1|1)$}{}{}
\end{pfloesung}
\begin{pfloesung}{}
\lz{25.}{Start $(0|3)$, 2 nach rechts, 1 nach oben: Gerade durch $(0|3)$ und $(2|4)$}{}{}
\end{pfloesung}
\begin{pfloesung}{}
\lz{26.}{Start $(0|-4)$, 1 nach rechts, 2 nach oben: Gerade durch $(0|-4)$ und $(1|-2)$}{}{}
\end{pfloesung}
\begin{pfloesung}{{\small\color{mbgrau}P10 2021 · 2a}}
\lz{27.}{Gerade durch (0; 1) und (1; 4)}{y-Achsenabschnitt 1, Steigungsdreieck 1 nach rechts, 3 nach oben}{}
\end{pfloesung}
\begin{pfloesung}{{\small\color{mbgrau}P10 2024 · 3a}}
\lz{28.}{Gerade durch (0|1) und (1|5), steil steigend}{Steigungsdreieck 1 nach rechts, 4 nach oben}{}
\end{pfloesung}
\begin{pfloesung}{{\small\color{mbgrau}P10 2026 FOR · 5a}}
\lz{29.}{Gerade durch (0|2) und (1|0), fallend}{Steigungsdreieck 1 nach rechts, 2 nach unten}{}
\end{pfloesung}
\begin{pfloesung}{{\small\color{mbgrau}P10 2022 · 3a}}
\lz{30.}{Gerade durch (0|3) und (1|1); x = 1,5}{0 = $-$2x + 3}{}
\end{pfloesung}
\begin{pfloesung}{}
\lz{31.}{$3x + 4$}{$m = (10 - 4) : (2 - 0) = 3$, $n = 4$: $f(x)$}{}
\end{pfloesung}
\begin{pfloesung}{}
\lz{32.}{$2x - 1$}{$m = (9 - 3) : (5 - 2) = 2$, $n = -1$: $f(x)$}{}
\end{pfloesung}
\begin{pfloesung}{}
\lz{33.}{$-3x + 9$}{$m = (0 - 6) : (3 - 1) = -3$, $n = 9$: $f(x)$}{}
\end{pfloesung}
\begin{pfloesung}{}
\lz{34.}{$2x - 3$}{$m = 2$, $n = -3$: $f(x)$}{}
\end{pfloesung}
\begin{pfloesung}{}
\lz{35.}{Gerade durch $A(-3|-1)$ und $B(2|4)$}{}{}
\end{pfloesung}
\begin{pfloesung}{{\small\color{mbgrau}P10 2017 · 5a}}
\lz{36.}{Gerade durch K und L; m = ½, n = 1 $\Rightarrow$ y = ½x + 1}{m = (2 $-$ ($-$1)) : (2 $-$ ($-$4)) = ½; 2 = ½ · 2 + n $\Rightarrow$ n = 1}{}
\end{pfloesung}
\begin{pfloesung}{{\small\color{mbgrau}P10 2015 · 4d}}
\lz{37.}{h(t) = $-$5t + 1300}{m = (700 $-$ 1200) : (120 $-$ 20) = $-$5; 1200 = $-$5 · 20 + n $\Rightarrow$ n = 1300}{}
\end{pfloesung}
\begin{pfloesung}{}
\lz{38.}{$4{,}20$: Grundgebühr 4{,}20 €}{$m = (25{,}20 - 12{,}60) : (10 - 4) = 2{,}10$; $n = 12{,}60 - 4 \cdot 2{,}10$}{}
\end{pfloesung}
\begin{pfloesung}{{\small\color{mbgrau}P10 2025 · 5a}}
\lz{39.}{Gerade durch A($-$2|6) und B(3|$-$1,5); Aussage 1 falsch (fallend), Aussage 2 wahr; f(x) = $-$1,5x + 3}{m = ($-$1,5 $-$ 6) : (3 + 2) = $-$1,5; n = 3}{}
\end{pfloesung}
\begin{pfloesung}{}
\lz{40.}{Gerade}{}{}
\end{pfloesung}
\begin{pfloesung}{}
\lz{41.}{Normalform}{}{}
\end{pfloesung}
\begin{pfloesung}{}
\lz{42.}{$13$}{$f(4) = 2 \cdot 4 + 5$}{}
\end{pfloesung}
\begin{pfloesung}{}
\lz{43.}{Schnittpunkt $(0|5)$}{$f(0) = 5$}{}
\end{pfloesung}
\begin{pfloesung}{}
\lz{44.}{$x_1 = 3$, $x_2 = -1$}{$0 = (x - 1)^2 - 4$, $4 = (x - 1)^2$, $x - 1 = 2$ oder $x - 1 = -2$}{}
\end{pfloesung}
\begin{pfloesung}{}
\lz{45.}{$36$}{$f(-6) = (-6)^2$}{}
\end{pfloesung}
\begin{pfloesung}{}
\lz{46.}{$-5$}{$f(-3) = 9 - 12 - 2$}{}
\end{pfloesung}
\begin{pfloesung}{{\small\color{mbgrau}P10 2017 · 5b}}
\lz{47.}{y = $-$4}{½ · ($-$10) + 1}{}
\end{pfloesung}
\begin{pfloesung}{{\small\color{mbgrau}P10 2026 FOR · 5b}}
\lz{48.}{Ja, f($-$4) = $-$2 · ($-$4) + 2 = 10}{f($-$4) = 10}{}
\end{pfloesung}
\begin{pfloesung}{}
\lz{49.}{$4$}{$4$; $1$; $0$; $1$}{}
\end{pfloesung}
\begin{pfloesung}{}
\lz{50.}{ja, $P$ liegt auf beiden Graphen}{$f(4) = 16 - 5 = 11$, $g(4) = 8 + 3 = 11$}{}
\end{pfloesung}
\begin{pfloesung}{}
\lz{51.}{Anfangswert 45 € , Änderung 38 € je Stunde}{$n = 45$; $m = 38$}{}
\end{pfloesung}
\begin{pfloesung}{}
\lz{52.}{3}{zum Beispiel $3 : 1 = 3$}{}
\end{pfloesung}
\begin{pfloesung}{}
\lz{53.}{$2x + 4$}{$y$}{}
\end{pfloesung}
\begin{pfloesung}{}
\lz{54.}{Gerade durch $(0|0)$, $(1|3)$, $(2|6)$}{y-Werte: $0$, $3$, $6$}{}
\end{pfloesung}
\begin{pfloesung}{}
\lz{55.}{$6x + 10$}{$m = 6$, $n = 10$: $y$}{}
\end{pfloesung}
\begin{pfloesung}{{\small\color{mbgrau}P10 2016 · 6c}}
\lz{56.}{y = 1,5x + 200}{}{}
\end{pfloesung}
\begin{pfloesung}{{\small\color{mbgrau}P10 2021 · 7a}}
\lz{57.}{Fahrpreis: y = 0,20x + 5; Fitnesscenter: y = 5x + 20}{20 ct = 0,20 €}{}
\end{pfloesung}
\begin{pfloesung}{{\small\color{mbgrau}P10 2021 · 6c}}
\lz{58.}{y: Höhe in cm; x: Zeit in min; 40: Anfangshöhe in cm}{}{}
\end{pfloesung}
\begin{pfloesung}{{\small\color{mbgrau}P10 2021 · 6d}}
\lz{59.}{200 min (3 h 20 min)}{$-$0,2x + 40 = 0 $\Rightarrow$ x = 200}{}
\end{pfloesung}
\begin{pfloesung}{{\small\color{mbgrau}P10 2022 · 6b}}
\lz{60.}{y = 55x + 990; 19 Monate}{55x + 990 $\ge$ 2000 $\Rightarrow$ x $\ge$ 18,36; aufrunden}{}
\end{pfloesung}
\begin{pfloesung}{}
\lz{61.}{x Minuten, y Liter}{zum Beispiel ein Tank mit 300 Litern, der je Minute 15 Liter verliert}{}
\end{pfloesung}
\begin{pfloesung}{}
\lz{62.}{$48x + 35$}{Nele hat Anfangswert und Änderung vertauscht; richtig: $m = 48$, $n = 35$, also $y$}{}
\end{pfloesung}
\begin{pfloesung}{}
\lz{63.}{$10$}{$3 \cdot 4 - 2$}{}
\end{pfloesung}
\begin{pfloesung}{}
\lz{64.}{$7$}{$5 \cdot 2 - 3$}{}
\end{pfloesung}
\begin{pfloesung}{}
\lz{65.}{$0$}{$m$}{}
\end{pfloesung}
\begin{pfloesung}{{\small\color{mbgrau}P10 2021 · 2b}}
\lz{66.}{Aussage 2 und 3 richtig}{f(5) = 16 $\neq$ 0}{}
\end{pfloesung}
\begin{pfloesung}{}
\lz{67.}{g, denn 5 ist größer als 3: je Schritt nach rechts geht g weiter nach oben}{}{}
\end{pfloesung}
\begin{pfloesung}{}
\lz{68.}{Nein: Beide haben dieselbe Steigung, aber verschiedene y-Achsenabschnitte}{}{}
\end{pfloesung}
\begin{pfloesung}{{\small\color{mbgrau}P10 2023 · 4a}}
\lz{69.}{Gerade durch (0|$-$2) und (1|2); S(1|2); Aussage 1 falsch, Aussage 2 wahr, Aussage 3 falsch}{Steigung 4, y-Achsenabschnitt $-$2; zweiter Schnittpunkt ($-$7|$-$30) außerhalb}{}
\end{pfloesung}
\begin{pfloesung}{}
\lz{70.}{Ja: $f(3) = 8$}{}{}
\end{pfloesung}
\begin{pfloesung}{}
\lz{71.}{$-4$}{$y = \tfrac{3}{4} \cdot (-8) + 2$}{}
\end{pfloesung}
\begin{pfloesung}{{\small\color{mbgrau}P10 2023 · 1i}}
\lz{72.}{(3|$-$24), denn f(3) = $-$18 $\neq$ $-$24}{f(3) = $-$18}{}
\end{pfloesung}
\begin{pfloesung}{}
\lz{73.}{$265$ €}{$85 + 12 \cdot 15$}{}
\end{pfloesung}
\begin{pfloesung}{{\small\color{mbgrau}P10 2021 · 6a}}
\lz{74.}{40 cm; 24 cm}{Abnahme 2 cm je 10 min}{}
\end{pfloesung}
\begin{pfloesung}{{\small\color{mbgrau}P10 2022 · 6a}}
\lz{75.}{1736 €}{12 · 22 = 264 €}{}
\end{pfloesung}
\begin{pfloesung}{}
\lz{76.}{Nein: $45 + 5 \cdot 52 = 305$ €, das sind 5 € zu viel}{}{}
\end{pfloesung}
\begin{pfloesung}{}
\lz{77.}{B}{Startwert auf der y-Achse und Anstieg passen}{}
\end{pfloesung}
\begin{pfloesung}{}
\lz{78.}{$1\,250$ €}{$830 + 12 \cdot 35$}{}
\end{pfloesung}
\begin{pfloesung}{{\small\color{mbgrau}P10 2016 · 6a}}
\lz{79.}{I Reiselust, II Sonnenschein; II beginnt bei 200 € (Grundpreis)}{Grundpreis = y-Achsenabschnitt = 200 €}{}
\end{pfloesung}
\begin{pfloesung}{}
\lz{80.}{A ist günstiger}{A: $50 + 12 \cdot 20 = 290$ €, B: $12 \cdot 28 = 336$ €}{}
\end{pfloesung}
\begin{pfloesung}{}
\lz{81.}{$60$ l}{zu Beginn $300$ l; nach 40 min: $300 - 8 \cdot 30$; 30 l mehr als nach 5 min}{}
\end{pfloesung}
\begin{pfloesung}{}
\lz{82.}{das ist nur dann der mit Grundgebühr, wenn er den kleineren Preis je km hat}{Nicht immer: Rechts vom Schnittpunkt ist der Tarif mit dem kleineren Preis je km günstiger}{}
\end{pfloesung}
\begin{pfloesung}{{\small\color{mbgrau}P10 2016 · 6b}}
\lz{83.}{Reiselust (700 € < 725 €); nein, bei 450 km Sonnenschein 875 € < 900 €}{200 + 1,5 · 350 = 725; 2 · 350 = 700}{}
\end{pfloesung}
\begin{pfloesung}{\textbf{Prüfstein} \textperiodcentered{} P10 2023 · 3}
\lz{a)}{Angebot 1: C; Angebot 2: B}{Angebot 1: konstant bis 350 km, dann steigend; Angebot 2: Grundgebühr, dann gleichmäßig steigend}{}
\lz{b)}{Angebot 1: 178,20 €; Angebot 2: 162,30 €; Angebot 2 günstiger; K(x) = 0,09x + 120}{120 Mehrkilometer · 0,36 € = 43,20 €; 5 · 27 = 135 €}{}
\end{pfloesung}
\end{document}